% =====================================================================
%  tikz-syntagmata.tex — έτοιμες συνταγές για τα οκτώ μοτίβα εξήγησης
%
%  Πρότυπο: τα σχήματα του άρθρου του Raft (Ongaro & Ousterhout, ATC '14).
%  Τα μοτίβα είναι εργαλεία σχεδίασης· κανένα σχήμα δεν αντιγράφεται.
%
%  ΧΡΗΣΗ: αντέγραψε το μπλοκ που θέλεις μέσα στο κεφάλαιο και άλλαξε τα
%  περιεχόμενα. Χρησιμοποιεί ΜΟΝΟ ό,τι φορτώνει ήδη το preamble.tex:
%  tikz (positioning, arrows.meta, calc, fit, backgrounds, shapes.geometric),
%  pgfplots, booktabs, xcolor με τα χρώματα cCFT cBFT cAux cGrey cSand cPurple.
%
%  ΟΛΑ ΤΑ ΜΠΛΟΚ ΕΧΟΥΝ ΜΕΤΑΓΛΩΤΤΙΣΤΕΙ ΜΕ XeLaTeX ΠΡΙΝ ΜΠΟΥΝ ΕΔΩ.
% =====================================================================


% ---------------------------------------------------------------------
%  Μ3 — ΔΙΑΓΡΑΜΜΑ ΚΑΤΑΣΤΑΣΕΩΝ            (πρότυπο: Raft, Σχ. 4)
%
%  Ο κανόνας: στα βέλη μπαίνει το ΓΕΓΟΝΟΣ που προκαλεί τη μετάβαση, σε
%  ελληνικά, όχι σύμβολο. Ο αναγνώστης διαβάζει το διάγραμμα σαν πρόταση.
% ---------------------------------------------------------------------
\begin{figure}[t]
\centering
\begin{tikzpicture}[
    font=\footnotesize,
    kat/.style={draw=cCFT,fill=cCFT!8,rounded corners=3pt,
                minimum height=7mm,minimum width=20mm,inner sep=3pt},
    >={Stealth[length=4pt]},
    every edge/.style={draw=cGrey,->,semithick}]

  \node[kat] (f) at (0,0)   {ακόλουθος};
  \node[kat] (y) at (4.2,0) {υποψήφιος};
  \node[kat] (h) at (8.4,0) {ηγέτης};

  \draw[->,semithick,cGrey] (f) -- node[above,align=center,font=\scriptsize]
        {λήξη χρονομέτρου·\\ ξεκινά εκλογή} (y);
  \draw[->,semithick,cGrey] (y) -- node[above,align=center,font=\scriptsize]
        {ψήφοι από\\ πλειοψηφία} (h);

  \draw[->,semithick,cGrey] (y) to[out=125,in=55,looseness=6]
        node[above,font=\scriptsize,align=center]
        {λήξη χρονομέτρου·\\ νέα εκλογή} (y);

  \draw[->,semithick,cGrey] (y) to[out=-145,in=-35]
        node[below,font=\scriptsize,align=center]
        {βρίσκει τρέχοντα ηγέτη\\ ή νεότερη εποχή} (f);
  \draw[->,semithick,cGrey] (h) to[out=-145,in=-35]
        node[below,font=\scriptsize,align=center]
        {βρίσκει κόμβο με\\ νεότερη εποχή} (y);

  \node[anchor=west,font=\scriptsize,text=cGrey] at (-2.6,0) {εκκίνηση};
  \draw[->,semithick,cGrey] (-1.3,0) -- (f);
\end{tikzpicture}
\caption[Καταστάσεις κόμβου και τι προκαλεί κάθε μετάβαση]%
{Οι τρεις καταστάσεις ενός κόμβου, και το γεγονός που προκαλεί κάθε
μετάβαση. \emph{Η λεζάντα οφείλει να λέει ολόκληρο τον μηχανισμό:} κάθε
κόμβος ξεκινά ως ακόλουθος και απαντά μόνο σε αιτήματα άλλων. Αν περάσει
το χρονόμετρο χωρίς να ακούσει ηγέτη, γίνεται υποψήφιος και ξεκινά εκλογή.
Ο υποψήφιος που συγκεντρώνει ψήφους από πλειοψηφία γίνεται ηγέτης, και
παραμένει ηγέτης ώσπου να αποτύχει ή να δει νεότερη εποχή.}
\label{fig:katastaseis}
\end{figure}


% ---------------------------------------------------------------------
%  Μ4 — ΓΡΑΜΜΗ ΧΡΟΝΟΥ ΜΕ ΕΠΟΧΕΣ         (πρότυπο: Raft, Σχ. 5)
%
%  Ο κανόνας: το ΣΠΑΣΜΕΝΟ σενάριο μπαίνει ΜΕΣΑ στο σχήμα, όχι σε δεύτερο
%  σχήμα. Εδώ η εποχή t3 τελειώνει χωρίς ηγέτη.
% ---------------------------------------------------------------------
\begin{figure}[t]
\centering
\begin{tikzpicture}[font=\footnotesize,x=1cm,y=1cm]
  \def\yb{0}
  % εποχές: x_αρχή / πλάτος_εκλογής / πλάτος_λειτουργίας / ετικέτα
  \foreach \x/\we/\wo/\lbl in {0/0.35/2.0/{εποχή 1}, 2.55/0.35/2.3/{εποχή 2},
                               5.4/0.35/0/{t3}, 6.0/0.35/2.0/{εποχή 4}}{
    \fill[cCFT]        (\x,\yb) rectangle ++(\we,0.42);
    \ifdim\wo pt>0pt
      \fill[cCFT!22] (\x+\we,\yb) rectangle ++(\wo,0.42);
    \fi
    \node[anchor=south,font=\scriptsize,text=cGrey]
         at (\x+\we/2+\wo/2,\yb+0.45) {\lbl};
  }
  \draw[->,semithick,cGrey] (-0.15,\yb) -- (8.4,\yb)
        node[anchor=west,font=\scriptsize] {χρόνος};

  \draw[cBFT,->,semithick] (0.18,-0.75) -- (0.18,-0.05);
  \node[anchor=north,font=\scriptsize,text=cBFT,align=center]
       at (0.18,-0.75) {εκλογή};
  \draw[cGrey,->,semithick] (1.6,-0.75) -- (1.6,-0.05);
  \node[anchor=north,font=\scriptsize,text=cGrey,align=center]
       at (1.6,-0.75) {κανονική\\λειτουργία};
  \draw[cBFT,->,semithick] (5.55,-0.75) -- (5.55,-0.05);
  \node[anchor=north,font=\scriptsize,text=cBFT,align=center]
       at (5.75,-0.75) {καμία εκλογή\\δεν πέτυχε};
\end{tikzpicture}
\caption[Ο χρόνος σε εποχές, με μία εποχή που τελειώνει χωρίς ηγέτη]%
{Ο χρόνος χωρίζεται σε \emph{εποχές} με αύξοντα αριθμό, και κάθε εποχή
ανοίγει με εκλογή. Μετά από πετυχημένη εκλογή, ένας ηγέτης διαχειρίζεται
τη συστάδα ως το τέλος της εποχής. \textbf{Η εποχή t3 δείχνει την
περίπτωση που πάει στραβά:} η ψήφος μοιράστηκε, κανένας δεν συγκέντρωσε
πλειοψηφία, και η εποχή τελείωσε χωρίς ηγέτη. Οι κόμβοι δεν βλέπουν
απαραίτητα τις μεταβάσεις την ίδια στιγμή.}
\label{fig:epoxes}
\end{figure}


% ---------------------------------------------------------------------
%  Μ5 — ΔΟΜΗ ΔΕΔΟΜΕΝΩΝ ΩΣ ΚΟΥΤΑΚΙΑ      (πρότυπο: Raft, Σχ. 6)
% ---------------------------------------------------------------------
\begin{figure}[t]
\centering
\begin{tikzpicture}[font=\scriptsize,x=0.92cm,y=0.72cm]
  \foreach \r/\name/\vals in {%
      0/{κόμβος 1}/{1,1,1,2,3},
      1/{κόμβος 2}/{1,1,1,2},
      2/{κόμβος 3}/{1,1,1,2,3}}{
    \node[anchor=east,font=\scriptsize] at (-0.25,-\r) {\name};
    \foreach \v [count=\c from 0] in \vals {
      \node[draw=cGrey,fill=white,minimum width=0.72cm,minimum height=0.52cm,
            inner sep=0pt] at (\c,-\r) {\v};
    }
  }
  \foreach \c in {0,1,2,3,4}{
    \node[font=\tiny,text=cGrey,anchor=south] at (\c,0.42) {\the\numexpr\c+1\relax};
  }
  \node[font=\tiny,text=cGrey,anchor=south] at (2,1.0) {θέση στο \en{log}};
  \draw[cCFT,thick] (-0.42,-2.42) rectangle (2.42,0.32);
  \node[anchor=north,font=\scriptsize,text=cCFT,align=center] at (1.0,-2.5)
       {δεσμευμένες: υπάρχουν σε πλειοψηφία};
\end{tikzpicture}
\caption[Το \en{log} ως αριθμημένα κουτάκια, και τι σημαίνει «δεσμευμένη»]%
{Κάθε κουτάκι είναι μία καταχώριση· ο αριθμός μέσα του είναι η εποχή στην
οποία δημιουργήθηκε. Οι τρεις πρώτες θέσεις υπάρχουν και στους τρεις
κόμβους, άρα σε πλειοψηφία, άρα είναι \emph{δεσμευμένες} και δεν
αναιρούνται ποτέ. Η θέση 4 υπάρχει σε δύο από τους τρεις — οριακά
δεσμευμένη. Η θέση 5 υπάρχει σε δύο κόμβους αλλά λείπει από τον κόμβο 2:
ακόμη μπορεί να χαθεί αν αλλάξει ηγέτης.}
\label{fig:log}
\end{figure}


% ---------------------------------------------------------------------
%  Μ2 — ΠΙΝΑΚΑΣ ΙΔΙΟΤΗΤΩΝ                (πρότυπο: Raft, Σχ. 3)
%
%  Ο κανόνας: ΟΝΟΜΑ ιδιότητας + ΜΙΑ πρόταση + δείκτης ενότητας. Το όνομα
%  είναι το σημαντικό: δίνει στον αναγνώστη λαβή για να την ξανασυζητήσει.
% ---------------------------------------------------------------------
\begin{table}[t]
\centering
\caption[Οι ιδιότητες που εγγυάται το πρωτόκολλο]%
{Οι ιδιότητες που το πρωτόκολλο εγγυάται \emph{ανά πάσα στιγμή}. Κάθε μία
έχει όνομα, ώστε τα επόμενα κεφάλαια να μπορούν να την επικαλεστούν χωρίς
να την ξαναδιατυπώσουν. Ο δείκτης δείχνει πού αποδεικνύεται.}
\label{tab:idiotites}
\small
\begin{tabular}{@{}l p{0.62\linewidth} l@{}}
\toprule
\textbf{Ιδιότητα} & \textbf{Τι εγγυάται} & \textbf{Πού} \\
\midrule
Μοναδικός ηγέτης   & Σε κάθε εποχή εκλέγεται το πολύ ένας ηγέτης.
                   & §5.2 \\
Μόνο προσθήκη      & Ο ηγέτης ποτέ δεν σβήνει ούτε αντικαθιστά καταχωρίσεις
                     του δικού του \en{log}· μόνο προσθέτει.  & §5.3 \\
Ταύτιση \en{log}   & Αν δύο \en{log} έχουν καταχώριση με ίδια θέση και ίδια
                     εποχή, τότε ταυτίζονται σε όλες τις προηγούμενες θέσεις.
                   & §5.3 \\
\bottomrule
\end{tabular}
\end{table}


% ---------------------------------------------------------------------
%  Μ6 — ΠΛΕΓΜΑ ΣΕΝΑΡΙΩΝ                  (πρότυπο: Raft, Σχ. 7)
%
%  Ο κανόνας: ΟΛΕΣ οι περιπτώσεις σε ΕΝΑ σχήμα, με γράμμα η καθεμία, και η
%  λεζάντα να λέει τι δείχνει κάθε γράμμα. Δύο σχήματα θα έκρυβαν ότι
%  πρόκειται για παραλλαγές του ίδιου πράγματος.
% ---------------------------------------------------------------------
\begin{figure}[t]
\centering
\begin{tikzpicture}[font=\scriptsize,x=0.55cm,y=0.62cm]
  \node[anchor=east,font=\scriptsize\bfseries] at (-0.4,0) {ηγέτης};
  \foreach \v [count=\c from 0] in {1,1,1,4,4,5,5,6,6,6}{
    \node[draw=cCFT,fill=cCFT!12,minimum width=0.46cm,minimum height=0.42cm,
          inner sep=0pt,font=\tiny] at (\c,0) {\v};}
  \foreach \r/\lbl/\vals in {%
      1/{(α)}/{1,1,1,4,4,5,5,6},
      2/{(β)}/{1,1,1,4},
      3/{(γ)}/{1,1,1,4,4,5,5,6,6,6,6},
      4/{(δ)}/{1,1,1,4,4,5,5,6,6,6,7,7},
      5/{(ε)}/{1,1,1,4,4,4,4},
      6/{(στ)}/{1,1,1,2,2,2,3,3,3,3,3}}{
    \node[anchor=east,font=\scriptsize] at (-0.4,-\r) {\lbl};
    \foreach \v [count=\c from 0] in \vals {
      \node[draw=cGrey,fill=white,minimum width=0.46cm,minimum height=0.42cm,
            inner sep=0pt,font=\tiny] at (\c,-\r) {\v};}
  }
\end{tikzpicture}
\caption[Έξι καταστάσεις που μπορεί να βρει ένας νέος ηγέτης]%
{Όταν ο ηγέτης της επάνω σειράς αναλαμβάνει, καθένας από τους ακολούθους
μπορεί να βρίσκεται σε οποιαδήποτε από τις καταστάσεις (α)–(στ). Κάθε
κουτάκι είναι μία καταχώριση και ο αριθμός μέσα του η εποχή της.
Στις (α) και (β) λείπουν καταχωρίσεις. Στις (γ) και (δ) υπάρχουν
\emph{περισσευούμενες} καταχωρίσεις από εποχές που δεν δεσμεύτηκαν ποτέ.
Οι (ε) και (στ) έχουν και τα δύο ταυτόχρονα. \textbf{Το σχήμα είναι ένα και
όχι έξι, επειδή ακριβώς αυτό θέλουμε να δει ο αναγνώστης: ότι πρόκειται για
παραλλαγές ενός προβλήματος, όχι για έξι διαφορετικά.}}
\label{fig:senaria}
\end{figure}
